\documentclass[runningheads]{llncs}
\usepackage{graphicx}
\usepackage{amsmath,amssymb}
\usepackage{booktabs}
\usepackage{color}
\usepackage[width=122mm,left=12mm,paperwidth=146mm,height=193mm,top=12mm,paperheight=217mm]{geometry}
\begin{document}
\pagestyle{headings}
\mainmatter
\def\ECCV16SubNumber{***}  % Insert your submission number here

\title{Hybridization vs. Co-Presence in Compositional Diffusion Models}

\titlerunning{ECCV-16 submission ID \ECCV16SubNumber}

\authorrunning{ECCV-16 submission ID \ECCV16SubNumber}

\author{Anonymous ECCV submission}
\institute{Paper ID \ECCV16SubNumber}

\maketitle

\begin{abstract}
We study the gap between semantic conjunction and logical composition
in Stable Diffusion 3.5.
We compare monolithic prompting (CLIP AND), SuperDiff logical composition,
and hybrid guidance under shared-noise trajectory analysis, and show that
the two operators follow structurally distinct latent paths that diverge
within the first few denoising steps.
To quantify this gap directly, we train an approximate closed-loop inverter
that maps SuperDiff-AND outputs back to the closest achievable SD3.5
conditioning vector $p^*$, enabling three-level residual measurement
across images, VAE latents, and full denoising trajectories.
We show that the residual gap after optimal inversion is predictable from
the CLIP geometry of the composed concept pair, providing an operational
taxonomy of composability.
Guided interpolation with moderate $\alpha$ narrows the trajectory gap
without collapsing to either extreme.
\keywords{compositional generation, diffusion models, latent trajectories,
approximate inversion, composability gap, logical operators}
\end{abstract}

% =========================================================================
\section{Introduction}
% =========================================================================

\textbf{Problem.}
\begin{itemize}
  \item Natural-language conjunction (``A and B'') and score-level logical
    composition via SuperDiff AND are distinct operators with different
    latent dynamics and different semantic outcomes.
  \item The product-of-experts distribution produced by SuperDiff AND is not
    in general representable by any single SD3.5 text prompt, but the
    expressibility gap is not uniform across concept pairs.
  \item Characterising this gap requires measuring it at the level of
    conditioning, latents, and trajectories --- not just decoded images.
\end{itemize}

\textbf{Objectives.}
\begin{itemize}
  \item Establish trajectory-level evidence that CLIP AND and SuperDiff AND
    are dynamically distinct under shared-noise conditions.
  \item Quantify the composability gap: how much of the SuperDiff-AND
    distribution can a single SD3.5 conditioning vector recover?
  \item Characterise what structural property of the concept pair predicts
    the gap magnitude.
  \item Evaluate hybrid guidance as a practical bridge between semantic and
    logical composition.
\end{itemize}

\textbf{Contributions.}
\begin{itemize}
  \item A shared-noise trajectory analysis protocol comparing monolithic,
    PoE, SuperDiff, and guided conditions with per-step divergence diagnostics
    and kappa tracking across 67 reproducible runs.
  \item An approximate inversion framework: a closed-loop CLIP+MLP inverter
    $f_\theta$ trained entirely within SD3.5 to find the best single
    conditioning vector $p^*$ approximating SuperDiff-AND outputs, without
    any external image-caption data.
  \item A three-level composability gap metric (CLIP cosine similarity and
    LPIPS at image level; MSE in VAE latent space; step-wise MSE and cosine
    similarity along matched denoising trajectories).
  \item A composability gap taxonomy that correlates residual gap magnitude
    with CLIP cosine similarity between the individual concept embeddings
    $e(c_1)$ and $e(c_2)$.
\end{itemize}

\textbf{Scope.}
\begin{itemize}
  \item Model: Stable Diffusion 3.5 Medium with FlowMatchEuler scheduler.
  \item Composition operators: monolithic CFG, SuperDiff AND (\texttt{fm\_ode}
    variant), PoE, guided hybrid, Composable NOT, SuperDiff NOT.
  \item All runs: fixed seed, 50 steps, guidance 4.5 unless stated otherwise.
\end{itemize}

% =========================================================================
\section{Background}
% =========================================================================

% Intentionally left for final version.

% =========================================================================
\section{Dataset Preparation}
% =========================================================================

\textbf{Trajectory experiment corpus.}
\begin{itemize}
  \item 67 reproducible runs across four intent groups: pairwise AND
    composition (31), guided $\alpha$ sweeps (9), solo stress prompts (29),
    negation (1), spatial extension (1), and multi-prompt dynamics (7).
  \item Each run stores a \texttt{summary.json} with divergence onset,
    endpoint L2, CLIP probe prediction, and kappa statistics, enabling
    one-to-one traceability from artifact to paper claim.
  \item All trajectory runs share the same scheduler configuration, seed,
    and noise initialisation within each comparison group.
\end{itemize}

\textbf{Inverter training corpus.}
\begin{itemize}
  \item 85 curated prompts across three tiers: single concepts (15),
    simple conjunctions (40), and attribute-rich variants (30).
  \item 8 images per prompt generated by SD3.5 (seeds 0--7, 50 steps,
    guidance 4.5), yielding 680 training images at $512\times512$.
  \item Ground-truth conditioning stored as \texttt{conditioning.pt} per
    prompt: \texttt{seq\_embeds} $(1, 410, 4096)$ and
    \texttt{pooled\_embeds} $(1, 2048)$ from SD3.5's triple text encoders.
  \item Prompt-level 90/10 train/validation split; the four held-out test
    pairs (cat/dog, person/umbrella, person/car, car/truck) are excluded
    from training entirely.
\end{itemize}

% =========================================================================
\section{Methodology}
% =========================================================================

\textbf{Part I\@: Shared-noise trajectory analysis.}
\begin{itemize}
  \item All conditions (prompt A alone, prompt B alone, monolithic ``A and B'',
    SuperDiff AND, PoE, guided hybrid) start from the same initial noise
    $x_T$ via a fixed generator seed.
  \item Latent trajectories recorded at every denoising step by
    \texttt{LatentTrajectoryCollector}, storing $(z_t, v_t, \sigma_t)$
    for all $t$.
  \item \textbf{Divergence onset}: first step at which pairwise L2 distance
    between any two conditions becomes non-trivial.
  \item \textbf{Endpoint L2}: final latent-space distance between conditions
    after all denoising steps.
  \item \textbf{Kappa diagnostics}: per-step concept weights $\kappa(t)$
    from the SuperDiff Proposition~6 solver, tracking how the composition
    rebalances over time.
  \item \textbf{CLIP probe}: nearest text prototype in CLIP space for each
    decoded output, used to label the effective semantic outcome.
  \item \textbf{Guided hybrid rule}: $v_\alpha = (1-\alpha)\,v_{\text{mono}}
    + \alpha\,v_{\text{superdiff}}$, swept over
    $\alpha \in \{0.0, 0.1, 0.3, 0.5, 0.7, 1.0\}$.
\end{itemize}

\textbf{Part II\@: Approximate inversion and composability gap measurement.}
\begin{itemize}
  \item \textbf{Inverter architecture} $f_\theta$: frozen CLIP ViT-L/14
    visual encoder (768-dim patch tokens, 257 positions including
    \texttt{[CLS]}); pooled head MLP
    $768 \!\to\! 1024 \!\to\! 2048$ predicts
    \texttt{pooled\_prompt\_embeds}; cross-attention sequence head
    (154 learned query tokens attending over patch tokens, projected
    $768 \!\to\! 4096$) predicts the CLIP portion of
    \texttt{prompt\_embeds}; T5 portion zeroed.
    Total trainable parameters: $\approx\!8\text{M}$.
  \item \textbf{Closed-loop training objective}:
    $\mathcal{L} = \mathrm{MSE}(\hat{e}_{\text{pool}}, e_{\text{pool}})
    + \lambda\,\mathrm{MSE}(\hat{e}_{\text{seq}}, e_{\text{seq}}[{:}154])$,
    where targets are the conditioning vectors SD3.5 itself produced
    when generating the training images.
    Trained with AdamW, cosine schedule, 50 epochs.
  \item \textbf{Test protocol}: for each held-out pair $(c_1, c_2)$,
    generate $N\!=\!8$ images from SuperDiff-AND; invert each image with
    $f_\theta$ to obtain $p^*$; regenerate from SD3.5($p^*$) from the
    same initial noise; compare to a naive monolithic baseline
    $c_1$ ``and'' $c_2$.
  \item \textbf{Three-level gap}:
    \emph{image} --- CLIP cosine similarity and LPIPS between AND and $p^*$
    outputs;
    \emph{latent} --- MSE in VAE encoder space at the final step;
    \emph{trajectory} --- step-wise MSE and cosine similarity between
    matched denoising paths under identical initial noise.
  \item \textbf{Characterisation}: CLIP cosine similarity between
    $e(c_1)$ and $e(c_2)$ (SD3.5 pooled conditioning vectors) used as
    the predictor variable; regressed against each gap metric to identify
    what geometric property of the concept pair determines composability.
\end{itemize}

% =========================================================================
\section{Prompts Tested}
% =========================================================================

\begin{itemize}
  \item \textbf{Object co-presence}: ``a dog'' + ``a cat''; verifies
    baseline logical AND behaviour.
  \item \textbf{Spatial relations}: ``a book on the left'' + ``a bird on
    the right''; ``cat on the left'' + ``dog on the right''; tests
    directional grounding under composition.
  \item \textbf{Attribute binding}: ``two blue apples'' + ``a red table'';
    ``a brown bench'' + ``a white building''; tests cross-object attribute
    preservation.
  \item \textbf{Negation}: ``a woman wearing a hat'' with NOT ``glasses'';
    ``a cake with strawberries'' with NOT ``chocolate''; compares Composable
    NOT (fixed weight) vs.\ SuperDiff NOT (dynamic $\kappa$).
  \item \textbf{Counting}: ``exactly three red apples''; ``three dogs and
    four cats''; tests cardinality control.
  \item \textbf{Multi-constraint (3-prompt)}: ``a cat and a dog'' +
    ``cat on the left'' + ``dog on the right''; tests semantic anchoring
    with simultaneous logical constraints.
  \item \textbf{Held-out inversion pairs}: cat/dog (complementary),
    person/umbrella (hierarchical), person/car (co-occurring), car/truck
    (same-category); chosen to span the CLIP cosine similarity spectrum.
\end{itemize}

% =========================================================================
\section{Semantic and Logical Composition Diverge in Latent Space}
% =========================================================================

\textbf{Hypothesis.}
\begin{itemize}
  \item If CLIP AND and SuperDiff AND were equivalent operators, they would
    follow the same latent trajectory from the same $x_T$.
    Any observed split is structural, not stochastic.
\end{itemize}

\textbf{Protocol.}
\begin{itemize}
  \item Conditions run under identical seed, steps, and guidance scale.
  \item Trajectory distances computed at every step; divergence onset
    recorded as the first step where L2 between CLIP and SuperDiff
    trajectories exceeds a noise-floor threshold.
  \item SuperDiff variant: \texttt{fm\_ode} with \texttt{--no-poe}.
\end{itemize}

\textbf{Findings.}
\begin{itemize}
  \item CLIP AND and SuperDiff AND diverge at steps 2--4 for all tested
    pairs; divergence is a property of the composition operator, not
    of late-stage refinement.
  \item SuperDiff endpoint geometry sits closer to a latent midpoint
    between the two single-concept endpoints than to the monolithic
    endpoint, reflecting the kappa balancing constraint.
  \item Kappa weights $\kappa(t)$ are time-varying with concept rebalancing
    concentrated in the first third of denoising, then stabilising.
  \item PoE and SuperDiff AND produce different trajectory shapes despite
    both being score-additive; the Proposition~6 normalisation is the
    decisive structural difference.
  \item Concrete: run \texttt{20260216\_114652} (cat-left/dog-right)
    shows divergence onset at step 3, endpoint
    L2(CLIP, SuperDiff) $= 377.7$.
    Run \texttt{20260217\_011236} (two-blue-apples/red-table)
    shows onset at step 2, L2 $= 407.1$.
\end{itemize}

% =========================================================================
\section{What Semantic Prompt Does Logical Composition Follow?}
% =========================================================================

\textbf{Hypothesis.}
\begin{itemize}
  \item The SuperDiff-AND distribution lies outside SD3.5's
    text-conditioning manifold for most concept pairs.
    An optimal single prompt $p^*$ can partially approximate it, but a
    residual gap will remain whose magnitude varies with the CLIP geometry
    of $(c_1, c_2)$.
  \item This residual --- measured after the best possible inversion ---
    is the operational definition of the composability gap.
\end{itemize}

\textbf{CLIP probe evidence (trajectory experiments).}
\begin{itemize}
  \item In run \texttt{20260217\_011236}, CLIP assigns the monolithic class
    to guided outputs but assigns a single sub-prompt class to pure
    SuperDiff, confirming that logical balancing and linguistic conjunction
    produce images in different semantic neighbourhoods.
  \item In run \texttt{20260216\_114652}, pure SuperDiff is pulled toward
    one component prompt under the CLIP probe, suggesting that when concepts
    are spatial, the AND operator can privilege one concept over the other.
\end{itemize}

\textbf{Inversion methodology.}
\begin{itemize}
  \item The inverter $f_\theta$ is calibrated within SD3.5's own loop;
    closed-loop sanity check (SD3.5($p$) $\to$ $f_\theta$ $\to$
    SD3.5($p^*$)) establishes an upper bound on achievable fidelity.
  \item The trajectory-level gap between AND and $p^*$ isolates the
    component of the composability gap that is structurally
    inexpressible via any single prompt, not attributable to inverter
    approximation error.
\end{itemize}

\textbf{Expected taxonomy from gap characterisation.}
\begin{itemize}
  \item \emph{Complementary pairs} (cat/dog, high CLIP similarity): AND
    distribution close to a natural conjunction; small gap; $p^*$ likely
    readable as a coherent natural-language prompt.
  \item \emph{Hierarchical pairs} (person/umbrella): AND distribution
    expressible as object+accessory relation; moderate gap.
  \item \emph{Co-occurring pairs} (person/car): moderate gap; gap driven
    by pose and spatial configuration variation.
  \item \emph{Same-category pairs} (car/truck): near-synonym in CLIP space;
    smallest expected gap; AND distribution close to either concept alone.
\end{itemize}

% =========================================================================
\section{Adding the Semantic Prompt as a Third Constraint}
% =========================================================================

\textbf{Hypothesis.}
\begin{itemize}
  \item Blending monolithic velocity $v_{\text{mono}}$ with SuperDiff
    velocity $v_{\text{superdiff}}$ at weight $\alpha$ defines a family
    of trajectories interpolating between semantic naturalness and logical
    balancing.
  \item There exists a regime of moderate $\alpha$ that reduces trajectory
    distance to the CLIP AND path without eliminating the compositional
    correction from SuperDiff.
\end{itemize}

\textbf{Method.}
\begin{itemize}
  \item Hybrid velocity: $v_\alpha = (1-\alpha)\,v_{\text{mono}} +
    \alpha\,v_{\text{superdiff}}$, computed at every denoising step.
  \item Sweep: $\alpha \in \{0.0, 0.1, 0.3, 0.5, 0.7, 1.0\}$; all other
    controls fixed.
\end{itemize}

\textbf{Findings.}
\begin{itemize}
  \item $\alpha = 0$ recovers monolithic CFG exactly;
    $\alpha = 1$ recovers pure SuperDiff exactly.
  \item At $\alpha = 0.3$: L2(CLIP, Guided) $= 146.0$ vs.\
    L2(CLIP, SuperDiff) $= 247.1$ for book-left/bird-right
    (\texttt{run 20260217\_005615}); a 41\% reduction.
  \item At $\alpha = 0.3$: L2(CLIP, Guided) $= 129.3$ vs.\
    L2(CLIP, SuperDiff) $= 407.1$ for two-blue-apples/red-table
    (\texttt{run 20260217\_011236}); a 68\% reduction.
  \item The trade-off is concave in $\alpha$: gains plateau beyond
    $\alpha = 0.5$ while semantic coherence degrades, making $\alpha = 0.3$
    the practical operating point for most pairs.
  \item Guided interpolation does not close the trajectory-level gap
    entirely, consistent with the structural inexpressibility established
    by the inversion experiment.
\end{itemize}

% =========================================================================
\section{Benchmark}
% =========================================================================

\textbf{Trajectory experiment protocol.}
\begin{itemize}
  \item Shared-noise pairwise comparisons; divergence onset and endpoint L2
    reported on the same latent trajectory basis.
  \item SuperDiff variant \texttt{fm\_ode}, \texttt{--no-poe}, 50 steps,
    guidance 4.5, seed 42 throughout.
  \item Ablations sweep: guidance scale $\in \{3.0, 4.5, 6.0, 7.5\}$;
    steps $\in \{30, 40, 50, 75, 100\}$; lift
    $\in \{0.0, 0.1, 0.3, 0.5\}$; 5 seeds for robustness;
    PCA vs.\ MDS projection comparison.
\end{itemize}

\textbf{Inversion experiment protocol.}
\begin{itemize}
  \item Four held-out pairs; 8 seeds each; 50 steps, guidance 4.5.
  \item Gap reported at all three levels per pair; characterisation
    regression run over CLIP cosine similarity of concept embeddings.
  \item Monolithic baseline (``$c_1$ and $c_2$'') included as a
    no-inversion reference point.
\end{itemize}

\textbf{Key quantitative results.}
\begin{itemize}
  \item Divergence onset: steps 2--4 for all tested pairs.
  \item Endpoint L2 (CLIP vs.\ SuperDiff): 247--407 latent units.
  \item Guided $\alpha = 0.3$ reduces L2 to CLIP path by 41--68\%.
  \item Negation: Composable NOT vs.\ SuperDiff NOT onset at step 1,
    endpoint L2 $= 325.2$ (\texttt{run 20260217\_013025}).
\end{itemize}

% =========================================================================
\section{Conclusion}
% =========================================================================

% Intentionally left for final version.

\bibliographystyle{splncs}
\bibliography{egbib}
\end{document}
